% Options for packages loaded elsewhere
\PassOptionsToPackage{unicode}{hyperref}
\PassOptionsToPackage{hyphens}{url}
\PassOptionsToPackage{dvipsnames,svgnames,x11names}{xcolor}
%
\documentclass[
  10pt,
  a4paper,
]{article}
\usepackage{amsmath,amssymb}
\usepackage{iftex}
\ifPDFTeX
  \usepackage[T1]{fontenc}
  \usepackage[utf8]{inputenc}
  \usepackage{textcomp} % provide euro and other symbols
\else % if luatex or xetex
  \usepackage{unicode-math} % this also loads fontspec
  \defaultfontfeatures{Scale=MatchLowercase}
  \defaultfontfeatures[\rmfamily]{Ligatures=TeX,Scale=1}
\fi
\usepackage{lmodern}
\ifPDFTeX\else
  % xetex/luatex font selection
\fi
% Use upquote if available, for straight quotes in verbatim environments
\IfFileExists{upquote.sty}{\usepackage{upquote}}{}
\IfFileExists{microtype.sty}{% use microtype if available
  \usepackage[]{microtype}
  \UseMicrotypeSet[protrusion]{basicmath} % disable protrusion for tt fonts
}{}
\makeatletter
\@ifundefined{KOMAClassName}{% if non-KOMA class
  \IfFileExists{parskip.sty}{%
    \usepackage{parskip}
  }{% else
    \setlength{\parindent}{0pt}
    \setlength{\parskip}{6pt plus 2pt minus 1pt}}
}{% if KOMA class
  \KOMAoptions{parskip=half}}
\makeatother
\usepackage{xcolor}
\usepackage[margin=23mm]{geometry}
\usepackage{color}
\usepackage{fancyvrb}
\newcommand{\VerbBar}{|}
\newcommand{\VERB}{\Verb[commandchars=\\\{\}]}
\DefineVerbatimEnvironment{Highlighting}{Verbatim}{commandchars=\\\{\}}
% Add ',fontsize=\small' for more characters per line
\newenvironment{Shaded}{}{}
\newcommand{\AlertTok}[1]{\textcolor[rgb]{1.00,0.00,0.00}{\textbf{#1}}}
\newcommand{\AnnotationTok}[1]{\textcolor[rgb]{0.38,0.63,0.69}{\textbf{\textit{#1}}}}
\newcommand{\AttributeTok}[1]{\textcolor[rgb]{0.49,0.56,0.16}{#1}}
\newcommand{\BaseNTok}[1]{\textcolor[rgb]{0.25,0.63,0.44}{#1}}
\newcommand{\BuiltInTok}[1]{\textcolor[rgb]{0.00,0.50,0.00}{#1}}
\newcommand{\CharTok}[1]{\textcolor[rgb]{0.25,0.44,0.63}{#1}}
\newcommand{\CommentTok}[1]{\textcolor[rgb]{0.38,0.63,0.69}{\textit{#1}}}
\newcommand{\CommentVarTok}[1]{\textcolor[rgb]{0.38,0.63,0.69}{\textbf{\textit{#1}}}}
\newcommand{\ConstantTok}[1]{\textcolor[rgb]{0.53,0.00,0.00}{#1}}
\newcommand{\ControlFlowTok}[1]{\textcolor[rgb]{0.00,0.44,0.13}{\textbf{#1}}}
\newcommand{\DataTypeTok}[1]{\textcolor[rgb]{0.56,0.13,0.00}{#1}}
\newcommand{\DecValTok}[1]{\textcolor[rgb]{0.25,0.63,0.44}{#1}}
\newcommand{\DocumentationTok}[1]{\textcolor[rgb]{0.73,0.13,0.13}{\textit{#1}}}
\newcommand{\ErrorTok}[1]{\textcolor[rgb]{1.00,0.00,0.00}{\textbf{#1}}}
\newcommand{\ExtensionTok}[1]{#1}
\newcommand{\FloatTok}[1]{\textcolor[rgb]{0.25,0.63,0.44}{#1}}
\newcommand{\FunctionTok}[1]{\textcolor[rgb]{0.02,0.16,0.49}{#1}}
\newcommand{\ImportTok}[1]{\textcolor[rgb]{0.00,0.50,0.00}{\textbf{#1}}}
\newcommand{\InformationTok}[1]{\textcolor[rgb]{0.38,0.63,0.69}{\textbf{\textit{#1}}}}
\newcommand{\KeywordTok}[1]{\textcolor[rgb]{0.00,0.44,0.13}{\textbf{#1}}}
\newcommand{\NormalTok}[1]{#1}
\newcommand{\OperatorTok}[1]{\textcolor[rgb]{0.40,0.40,0.40}{#1}}
\newcommand{\OtherTok}[1]{\textcolor[rgb]{0.00,0.44,0.13}{#1}}
\newcommand{\PreprocessorTok}[1]{\textcolor[rgb]{0.74,0.48,0.00}{#1}}
\newcommand{\RegionMarkerTok}[1]{#1}
\newcommand{\SpecialCharTok}[1]{\textcolor[rgb]{0.25,0.44,0.63}{#1}}
\newcommand{\SpecialStringTok}[1]{\textcolor[rgb]{0.73,0.40,0.53}{#1}}
\newcommand{\StringTok}[1]{\textcolor[rgb]{0.25,0.44,0.63}{#1}}
\newcommand{\VariableTok}[1]{\textcolor[rgb]{0.10,0.09,0.49}{#1}}
\newcommand{\VerbatimStringTok}[1]{\textcolor[rgb]{0.25,0.44,0.63}{#1}}
\newcommand{\WarningTok}[1]{\textcolor[rgb]{0.38,0.63,0.69}{\textbf{\textit{#1}}}}
\setlength{\emergencystretch}{3em} % prevent overfull lines
\providecommand{\tightlist}{%
  \setlength{\itemsep}{0pt}\setlength{\parskip}{0pt}}
\setcounter{secnumdepth}{-\maxdimen} % remove section numbering
\ifLuaTeX
  \usepackage{selnolig}  % disable illegal ligatures
\fi
\usepackage{bookmark}
\IfFileExists{xurl.sty}{\usepackage{xurl}}{} % add URL line breaks if available
\urlstyle{same}
\hypersetup{
  colorlinks=true,
  linkcolor={black},
  filecolor={Maroon},
  citecolor={Blue},
  urlcolor={blue},
  pdfcreator={LaTeX via pandoc}}

\author{}
\date{}

\begin{document}

\section{DDTA BA5 canonical semantic registry and controlled-authoring
candidate -
R2}\label{ddta-ba5-canonical-semantic-registry-and-controlled-authoring-candidate---r2}

\textbf{Status:} PROVISIONAL CANDIDATE AFTER BA5-T2 / BA5 OPEN

\textbf{Repository baseline reviewed:}
\texttt{85622414b2ff52d58f3cd11776fef4b3753afc7d}

\textbf{Supersedes for active BA5 reasoning:} the R1 canonical-referent
naming candidate. The R1 artifact remains immutable BA5-T1 derivation
history.

\textbf{Closed dependencies:} documentation layer; BA0 responsibility
boundary; BA1 \texttt{BAReferent\ +\ BAProposition}; BA2
proposition/operator/role/classification contract; BA3
provenance/continuity; BA4 projection contract; BA5-T1 canonical
referent naming lower bound.

\subsection{1. Purpose}\label{purpose}

BA5-T2 asks whether the strict T1 authoring rule can cover the wider
operative Base Analysis vocabulary without introducing a synonym
ontology or letting a tool decide semantic equivalence.

The starting rule remains:

\begin{Shaded}
\begin{Highlighting}[]
\NormalTok{registered canonical token}
\NormalTok{    {-}\textgreater{} normative semantic authoring}

\NormalTok{unregistered alias / synonym}
\NormalTok{    {-}\textgreater{} not silently equivalent}
\end{Highlighting}
\end{Shaded}

T2 finds that the rule survives, but a single flat registry does not.
Canonical tokens belong to semantic domains with different scoping,
authority and extension rules.

\subsection{2. Provisional R2 lower
bound}\label{provisional-r2-lower-bound}

The smallest current-scope registry model is a logical set of
domain-scoped registries, not one universal token table:

\begin{Shaded}
\begin{Highlighting}[]
\NormalTok{CanonicalSemanticRegistrySet                    [logical umbrella; NOT BAE]}
\NormalTok{|}
\NormalTok{+{-}{-} CanonicalReferentNameRegistry                [project{-}governed]}
\NormalTok{|    |{-} governedBaselineKey}
\NormalTok{|    |{-} namingScopeKey}
\NormalTok{|    \textasciigrave{}{-} canonicalName {-}\textgreater{} governed referent binding}
\NormalTok{|}
\NormalTok{+{-}{-} SemanticOperatorRegistry                     [BA2 contract{-}governed]}
\NormalTok{|    |{-} registryRevisionKey immutable}
\NormalTok{|    \textasciigrave{}{-} semanticOperatorKey {-}\textgreater{} normative BA2 meaning}
\NormalTok{|}
\NormalTok{+{-}{-} OperatorRoleRegistry                         [BA2 contract{-}governed]}
\NormalTok{|    |{-} registryRevisionKey immutable}
\NormalTok{|    \textasciigrave{}{-} (semanticOperatorKey, roleKey) {-}\textgreater{} BA2 role contract}
\NormalTok{|}
\NormalTok{+{-}{-} SemanticKindRegistry                         [method{-}neutral, extensible]}
\NormalTok{|    |{-} registryRevisionKey immutable}
\NormalTok{|    \textasciigrave{}{-} semanticKindKey {-}\textgreater{} normative method{-}neutral definition}
\NormalTok{|}
\NormalTok{\textasciigrave{}{-}{-} ControlledValueRegistry 0..*                 [domain{-}scoped, extensible where allowed]}
\NormalTok{     |{-} valueDomainKey}
\NormalTok{     |{-} registryRevisionKey immutable}
\NormalTok{     \textasciigrave{}{-} canonicalValueKey {-}\textgreater{} normative domain meaning}
\end{Highlighting}
\end{Shaded}

This is a semantic responsibility model. It does not require five
databases, five files, graph node kinds or ThreatForge classes. Physical
co-location is allowed when domain, scope, authority and revision remain
mechanically recoverable.

\subsection{3. Why one flat namespace is
rejected}\label{why-one-flat-namespace-is-rejected}

Token equality across domains does not establish semantic equality.

For example, a project might legitimately contain a referent named
\texttt{source} while BA2 also uses \texttt{source} as the role of
\texttt{transfer}.

\begin{Shaded}
\begin{Highlighting}[]
\NormalTok{referent canonicalName = source}
\NormalTok{roleKey under transfer = source}
\end{Highlighting}
\end{Shaded}

Those tokens occupy different domains and must not collide or alias
merely because their strings are identical.

Therefore uniqueness is domain-scoped:

\begin{Shaded}
\begin{Highlighting}[]
\NormalTok{same domain + same active scope + same canonical token}
\NormalTok{    {-}\textgreater{} one registry meaning/binding}

\NormalTok{same token in different domains}
\NormalTok{    {-}\textgreater{} allowed when each domain resolution is unambiguous}
\end{Highlighting}
\end{Shaded}

A global token uniqueness rule is \textbf{REJECTED AS OVER-STRONG}.

\subsection{4. Referent names retain the T1
rule}\label{referent-names-retain-the-t1-rule}

For a named shared project referent:

\begin{Shaded}
\begin{Highlighting}[]
\NormalTok{one governed baseline / naming scope}
\NormalTok{    {-}\textgreater{} one exact canonicalName}
\end{Highlighting}
\end{Shaded}

The canonical name is project-governed lexical metadata, not BA
identity. A governed rename may retain \texttt{BAReferent} identity
through BA3 continuity.

Nothing in T2 weakens the T1 requirement that shared views use the same
canonical referent name when presenting that shared entity.

\subsection{5. Semantic operator registry
coverage}\label{semantic-operator-registry-coverage}

All thirteen BA2 operator keys pass exact-token enforcement:

\begin{Shaded}
\begin{Highlighting}[]
\NormalTok{transfer}
\NormalTok{produce}
\NormalTok{create}
\NormalTok{observe}
\NormalTok{transition}
\NormalTok{correlate}
\NormalTok{reference}
\NormalTok{dependOn}
\NormalTok{consumeService}
\NormalTok{realize}
\NormalTok{assignResponsibility}
\NormalTok{constrain}
\NormalTok{classify}
\end{Highlighting}
\end{Shaded}

These tokens are not project-local synonyms. Their normative meaning
belongs to the accepted BA2 methodology contract.

Examples:

\begin{Shaded}
\begin{Highlighting}[]
\NormalTok{send            {-}\textgreater{} not an accepted semanticOperatorKey}
\NormalTok{transfer        {-}\textgreater{} accepted canonical key}

\NormalTok{usesService     {-}\textgreater{} not an accepted semanticOperatorKey}
\NormalTok{consumeService  {-}\textgreater{} accepted canonical key}
\end{Highlighting}
\end{Shaded}

Free explanatory prose may use natural language, but a
machine-significant operator binding must resolve to the exact
registered key. Hidden lexical mapping from \texttt{send} to
\texttt{transfer} is not part of the BA5 core.

The operator registry is scoped to an immutable methodology/BA2 registry
revision, not separately renamed by every project baseline.

\subsection{6. Operator-scoped role
registry}\label{operator-scoped-role-registry}

BA2 already rejects a context-free global role contract. T2 therefore
rejects a flat role registry that stores \texttt{roleKey} without its
operator context.

The canonical lookup key is conceptually:

\begin{Shaded}
\begin{Highlighting}[]
\NormalTok{(semanticOperatorKey, roleKey)}
\end{Highlighting}
\end{Shaded}

Examples from the closed BA2 contract:

\begin{Shaded}
\begin{Highlighting}[]
\NormalTok{transfer:}
\NormalTok{  source}
\NormalTok{  destination}
\NormalTok{  content}

\NormalTok{produce:}
\NormalTok{  actor}
\NormalTok{  result}
\NormalTok{  input}

\NormalTok{create:}
\NormalTok{  actor}
\NormalTok{  created}

\NormalTok{observe:}
\NormalTok{  actor}
\NormalTok{  observed}
\NormalTok{  result}

\NormalTok{transition:}
\NormalTok{  subject}
\NormalTok{  toState}
\NormalTok{  actor}
\NormalTok{  fromState}

\NormalTok{correlate:}
\NormalTok{  correlatedItem}
\NormalTok{  correlationContext}

\NormalTok{reference:}
\NormalTok{  referencer}
\NormalTok{  referenced}

\NormalTok{dependOn:}
\NormalTok{  dependent}
\NormalTok{  prerequisite}

\NormalTok{consumeService:}
\NormalTok{  consumer}
\NormalTok{  service}
\NormalTok{  provider}

\NormalTok{realize:}
\NormalTok{  abstract}
\NormalTok{  realization}

\NormalTok{assignResponsibility:}
\NormalTok{  responsibleParty}
\NormalTok{  responsibilityScope}
\NormalTok{  responsibilityKind}

\NormalTok{constrain:}
\NormalTok{  constraintTarget}
\NormalTok{  constraintValue}

\NormalTok{classify:}
\NormalTok{  classifiedReferent}
\NormalTok{  semanticKind}
\end{Highlighting}
\end{Shaded}

The same lexical role token may recur under more than one operator
without implying a context-free contract. Cardinality and
required/optional semantics remain BA2 responsibilities; BA5 requires
exact contextual lookup rather than redefining them.

Alias attempts such as \texttt{src}, \texttt{dest}, \texttt{owner} or
\texttt{kind} do not silently normalize to registered roles.

\subsection{7. Semantic-kind registry}\label{semantic-kind-registry}

BA2 closes a registry contract but intentionally does not close a
universal semantic-kind taxonomy.

BA5-T2 therefore separates two questions:

\begin{Shaded}
\begin{Highlighting}[]
\NormalTok{Is semanticKind a controlled canonical domain?   YES}
\NormalTok{Is every possible semanticKind pre{-}enumerated?   NO}
\end{Highlighting}
\end{Shaded}

Current evidence gives clear exact examples such as:

\begin{Shaded}
\begin{Highlighting}[]
\NormalTok{store}
\NormalTok{contract}
\end{Highlighting}
\end{Shaded}

These can be used as canonical \texttt{semanticKind} tokens when their
normative method-neutral definitions are registered.

BA2 also discusses slash-form conceptual examples:

\begin{Shaded}
\begin{Highlighting}[]
\NormalTok{party/person}
\NormalTok{information/resource}
\NormalTok{behavior/event}
\end{Highlighting}
\end{Shaded}

BA5-T2 does \textbf{not} silently turn this explanatory wording into
canonical keys. Exact token admission must be explicit.

A useful-looking term such as:

\begin{Shaded}
\begin{Highlighting}[]
\NormalTok{channel}
\end{Highlighting}
\end{Shaded}

is therefore \textbf{CANDIDATE / EVIDENCE-GATED}, not automatically
accepted merely because \texttt{Interface\ /\ Connection\ /\ Channel}
was discussed during BA1 ontology pressure. If a governed corpus needs a
method-neutral \texttt{channel} distinction that existing kinds cannot
preserve, it may enter the governed extension workflow below.

\subsection{8. Controlled-value registry
domains}\label{controlled-value-registry-domains}

Not every non-referent value should become a registry token. T2
preserves the BA1/BA2 distinction between controlled vocabulary and
ordinary typed/local values.

A controlled-value registry is justified when a value is selected from a
reusable semantic domain whose distinctions must be stable across
authoring and consumers.

Current BA2 examples include:

\begin{Shaded}
\begin{Highlighting}[]
\NormalTok{polarity:}
\NormalTok{  affirmative}
\NormalTok{  negative}

\NormalTok{scopedModifierKind:}
\NormalTok{  condition}
\NormalTok{  temporalScope}

\NormalTok{responsibilityKind:}
\NormalTok{  ownership}
\NormalTok{  management}
\NormalTok{  ... evidence{-}gated additional authority modes}
\end{Highlighting}
\end{Shaded}

For \texttt{polarity} and \texttt{scopedModifierKind}, the accepted BA2
contract already closes the exact current-scope values.

For \texttt{responsibilityKind}, BA2 requires a controlled domain but
does not claim an exhaustive universal list. \texttt{ownership} and
\texttt{management} are current evidence-backed values; new authority
modes are admitted only through governance.

A literal such as a duration, number, URI, threshold or project-specific
string does not become a canonical registry term merely because it
appears in a proposition. If it needs independent reusable identity, BA1
promotes it to a \texttt{BAReferent}; otherwise it remains an
appropriate typed/local value.

\subsection{9. Registry revision
identity}\label{registry-revision-identity}

T2 finds that mutable registry meaning would break reproducibility even
when the canonical token string stays unchanged.

Therefore any registry whose definitions/contracts can evolve requires
immutable revision identity:

\begin{Shaded}
\begin{Highlighting}[]
\NormalTok{same token}
\NormalTok{+ registry revision R1}
\NormalTok{    may not be assumed equivalent to}
\NormalTok{same token}
\NormalTok{+ registry revision R2}
\end{Highlighting}
\end{Shaded}

A registry revision may preserve all prior meanings while adding a new
entry, but the used revision remains inspectable.

The referent-name registry is already anchored by an immutable governed
documentation baseline/naming scope. BA2 contract registries and
extensible semantic/value registries require equivalent immutable
revision resolution.

This refines registry governance; it does not create a new BAE
lifecycle. BA3 already assumes reproducibility under the same accepted
semantic registries.

\subsection{10. Exact-token path across documentation, BA and
projections}\label{exact-token-path-across-documentation-ba-and-projections}

The accepted path is:

\begin{Shaded}
\begin{Highlighting}[]
\NormalTok{governed semantic authoring}
\NormalTok{  {-}\textgreater{} exact domain{-}scoped canonical token}
\NormalTok{      {-}\textgreater{} accepted BA binding}
\NormalTok{          {-}\textgreater{} projection trace / shared rendering}
\end{Highlighting}
\end{Shaded}

For example:

\begin{Shaded}
\begin{Highlighting}[]
\NormalTok{referent: cameraIngresso}
\NormalTok{operator: transfer}
\NormalTok{role: transfer/source}
\NormalTok{kind: store}
\end{Highlighting}
\end{Shaded}

A human or method projection may use its own presentation/type
vocabulary, but it must not rewrite the underlying shared key when
claiming to expose shared BA semantics.

Allowed:

\begin{Shaded}
\begin{Highlighting}[]
\NormalTok{[FlowParticipant] cameraIngresso}
\NormalTok{method{-}owned kind = DataFlowEndpoint}
\NormalTok{BA trace {-}\textgreater{} transfer/source {-}\textgreater{} cameraIngresso}
\end{Highlighting}
\end{Shaded}

Rejected:

\begin{Shaded}
\begin{Highlighting}[]
\NormalTok{method label "sender"}
\NormalTok{  silently substitutes for BA roleKey source}
\end{Highlighting}
\end{Shaded}

Method-specific taxonomy remains projection-owned under BA4; it is not
admitted into the BA registry merely because it is useful.

\subsection{11. Governed extension
workflow}\label{governed-extension-workflow}

When authoring needs an unregistered semantic token, the tool or author
does not immediately add it.

The minimum decision path is:

\begin{Shaded}
\begin{Highlighting}[]
\NormalTok{unregistered semantic need}
\NormalTok{    |}
\NormalTok{    +{-}{-} lexical alias of existing registered meaning?}
\NormalTok{    |      {-}\textgreater{} REJECT new entry}
\NormalTok{    |      {-}\textgreater{} use existing canonical token}
\NormalTok{    |}
\NormalTok{    +{-}{-} method{-}specific/projection{-}local meaning?}
\NormalTok{    |      {-}\textgreater{} keep downstream in projection/method registry}
\NormalTok{    |}
\NormalTok{    +{-}{-} independently reusable project{-}semantic identity?}
\NormalTok{    |      {-}\textgreater{} represent as BAReferent, not controlled local value}
\NormalTok{    |}
\NormalTok{    +{-}{-} new entry inside an already extensible method{-}neutral domain?}
\NormalTok{    |      {-}\textgreater{} candidate registry extension}
\NormalTok{    |      {-}\textgreater{} evidence + normative definition + non{-}equivalence check}
\NormalTok{    |      {-}\textgreater{} governed review}
\NormalTok{    |      {-}\textgreater{} new immutable registry revision if accepted}
\NormalTok{    |}
\NormalTok{    \textasciigrave{}{-}{-} changes a closed BA2 operator/role/polarity/modifier contract?}
\NormalTok{           {-}\textgreater{} smallest BA2 reopen candidate}
\NormalTok{           {-}\textgreater{} BA5 cannot smuggle it in as lexical metadata}
\end{Highlighting}
\end{Shaded}

This is the central T2 authority boundary.

\subsection{12. Minimum extension
packet}\label{minimum-extension-packet}

For an extensible \texttt{semanticKind} or controlled-value domain, an
accepted new canonical entry requires at least:

\begin{enumerate}
\def\labelenumi{\arabic{enumi}.}
\tightlist
\item
  \texttt{registryDomain} / \texttt{valueDomainKey};
\item
  proposed exact canonical token;
\item
  normative methodology-neutral definition;
\item
  concrete governed evidence or integration need;
\item
  comparison against existing entries showing that it is not merely an
  alias;
\item
  confirmation that the meaning is not method-specific;
\item
  confirmation that independent BAE identity is not the actual missing
  requirement;
\item
  governed acceptance; and
\item
  publication in an immutable registry revision used by subsequent
  authoring/BA baselines.
\end{enumerate}

The tool may assemble and validate this packet. It may not approve it as
semantic authority.

\subsection{13. Closed BA2 contract versus extensible registry
content}\label{closed-ba2-contract-versus-extensible-registry-content}

T2 makes an important distinction:

\begin{Shaded}
\begin{Highlighting}[]
\NormalTok{new semanticKind / controlled value}
\NormalTok{within an already open registry contract}
\NormalTok{    {-}\textgreater{} may be admitted by governed registry extension}

\NormalTok{new semanticOperatorKey}
\NormalTok{or changed operator{-}scoped role contract}
\NormalTok{or new polarity/modifier kind}
\NormalTok{    {-}\textgreater{} changes closed BA2 semantics}
\NormalTok{    {-}\textgreater{} smallest BA2 reopen required}
\end{Highlighting}
\end{Shaded}

This prevents BA5 from becoming a back door for modifying BA2.

A hypothetical new \texttt{publish} operator that is merely conveyance
is rejected as an alias of \texttt{transfer}. If concrete evidence
proves a distinct method-neutral action not representable by the
thirteen accepted operators, that evidence belongs to BA2 reopen review,
not an automatic vocabulary extension.

\subsection{14. Tool-support boundary}\label{tool-support-boundary}

A future tool may:

\begin{itemize}
\tightlist
\item
  validate exact tokens against the correct domain/revision;
\item
  offer context-sensitive completion;
\item
  after selecting \texttt{transfer}, offer only its accepted role keys;
\item
  flag duplicate/colliding referent names inside the active naming
  scope;
\item
  flag unregistered semantic kinds or controlled values;
\item
  show normative definitions and registry provenance;
\item
  classify an unregistered token as a candidate extension request;
\item
  assemble evidence for human/governed review.
\end{itemize}

It may not:

\begin{itemize}
\tightlist
\item
  infer hidden synonym equivalence;
\item
  pick a new canonical token and treat it as accepted;
\item
  change an operator/role contract through a local project registry;
\item
  promote method taxonomy into BA vocabulary;
\item
  decide that a local value needs BAE identity; or
\item
  mutate a registry definition in place.
\end{itemize}

ThreatForge remains an enforcement/integration instrument, never the
semantic authority.

\subsection{15. T2 mutation pressure}\label{t2-mutation-pressure}

\subsubsection{V1 - operator alias}\label{v1---operator-alias}

\begin{Shaded}
\begin{Highlighting}[]
\NormalTok{send {-}\textgreater{} transfer}
\end{Highlighting}
\end{Shaded}

\texttt{send} is not accepted as a second normative operator token. Use
\texttt{transfer} in the semantic binding.

Result: \textbf{PASS / ALIAS REJECTED}.

\subsubsection{V2 - role alias and context
loss}\label{v2---role-alias-and-context-loss}

\begin{Shaded}
\begin{Highlighting}[]
\NormalTok{transfer/src}
\end{Highlighting}
\end{Shaded}

\texttt{src} is rejected. \texttt{source} is valid only through the
\texttt{transfer} role contract. A context-free role lookup is
insufficient.

Result: \textbf{PASS WITH DOMAIN/CONTEXT REFINEMENT}.

\subsubsection{V3 - same token in different
domains}\label{v3---same-token-in-different-domains}

\begin{Shaded}
\begin{Highlighting}[]
\NormalTok{referent canonicalName = source}
\NormalTok{transfer roleKey       = source}
\end{Highlighting}
\end{Shaded}

No collision exists because the registries are domain-scoped.

Result: \textbf{PASS; FLAT GLOBAL NAMESPACE REJECTED}.

\subsubsection{V4 - semantic-kind alias}\label{v4---semantic-kind-alias}

\begin{Shaded}
\begin{Highlighting}[]
\NormalTok{repository}
\end{Highlighting}
\end{Shaded}

If the proposed definition is materially the same as registered
\texttt{store}, reject the new token as a synonym. If governed evidence
shows a distinct method-neutral concept, process it as a new
semantic-kind candidate rather than guessing from the word.

Result: \textbf{PASS}.

\subsubsection{\texorpdfstring{V5 - candidate
\texttt{channel}}{V5 - candidate channel}}\label{v5---candidate-channel}

\texttt{channel} is not admitted merely because it is a familiar systems
term. No T2 evidence proves that current
\texttt{BAReferent\ +\ classify\ +\ transfer/realize} semantics lose
required shared meaning without it.

Result: \textbf{DEFERRED / EVIDENCE-GATED}.

\subsubsection{V6 - responsibility kind}\label{v6---responsibility-kind}

\texttt{ownership} and \texttt{management} remain controlled values. A
proposed \texttt{control} value must receive a normative distinction and
evidence; it cannot enter as a loose synonym.

Result: \textbf{PASS}.

\subsubsection{V7 - genuinely new
operator}\label{v7---genuinely-new-operator}

A proposed operator whose meaning is not representable by the
thirteen-key BA2 registry cannot be accepted through BA5 alone.

Result: \textbf{BA2 REOPEN PATH REQUIRED IF MATERIAL EVIDENCE SURVIVES};
no current reopen triggered.

\subsubsection{V8 - projection
terminology}\label{v8---projection-terminology}

A method may display \texttt{DataFlowEndpoint}, \texttt{TrustBoundary}
or another method-owned type while tracing to canonical BA keys. Those
labels do not enter the BA registry.

Result: \textbf{PASS}.

\subsubsection{V9 - tool autocomplete}\label{v9---tool-autocomplete}

With \texttt{semanticOperatorKey\ =\ transfer}, the tool may offer
\texttt{source}, \texttt{destination}, \texttt{content} from the
immutable role contract. It may not offer \texttt{src} or invent another
role based on usage frequency.

Result: \textbf{PASS}.

\subsection{16. Reopen checks}\label{reopen-checks}

No current T2 case forces a reopen of BA0-BA4.

\begin{itemize}
\tightlist
\item
  BA1: registry entries are not a third semantic identity family;
  independent identity still promotes to \texttt{BAReferent}.
\item
  BA2: current operators/roles remain sufficient for the tested cases.
  The extension workflow explicitly routes a genuinely new operator/role
  back to BA2 rather than changing it silently.
\item
  BA3: immutable registry revision resolution is compatible with its
  existing reproducibility assumption and does not add a project
  lifecycle.
\item
  BA4: method-local taxonomy remains downstream; canonical BA token
  trace remains sufficient.
\item
  Documentation metamodel: controlled authoring remains a cross-document
  governance constraint, not a new MR/Decision/FR/SR family.
\end{itemize}

\subsection{17. Provisional
dispositions}\label{provisional-dispositions}

\begin{Shaded}
\begin{Highlighting}[]
\NormalTok{strict canonical{-}token authoring                         REQUIRED}
\NormalTok{one flat global token namespace                          REJECTED}
\NormalTok{registry domain/context resolution                       REQUIRED}
\NormalTok{referent canonicalName baseline/naming scoped            REQUIRED}
\NormalTok{13 BA2 operator keys exact                               REQUIRED}
\NormalTok{operator role lookup as (operatorKey, roleKey)           REQUIRED}
\NormalTok{role alias/synonym normalization                          REJECTED}
\NormalTok{semanticKind as controlled extensible domain              REQUIRED}
\NormalTok{universal semanticKind taxonomy                           REJECTED / NOT FORCED}
\NormalTok{store / contract as current exact semanticKind evidence  ACCEPTED FOR T2 PRESSURE}
\NormalTok{channel                                                   DEFERRED / EVIDENCE{-}GATED}
\NormalTok{controlled reusable value domains                         REQUIRED WHERE APPLICABLE}
\NormalTok{literal/local value {-}\textgreater{} registry entry by default          REJECTED}
\NormalTok{polarity affirmative|negative                             ACCEPTED EXACT KEYS}
\NormalTok{modifier kind condition|temporalScope                     ACCEPTED EXACT KEYS}
\NormalTok{responsibilityKind controlled domain                      REQUIRED}
\NormalTok{ownership / management current values                     ACCEPTED FOR T2 PRESSURE}
\NormalTok{immutable registry revision resolution                    REQUIRED}
\NormalTok{alias candidate {-}\textgreater{} new registry entry                     REJECTED}
\NormalTok{method{-}specific term {-}\textgreater{} BA registry                       REJECTED}
\NormalTok{governed extensible{-}domain admission                      REQUIRED}
\NormalTok{new operator/role through BA5 extension                   REJECTED}
\NormalTok{material new operator/role {-}\textgreater{} smallest BA2 reopen         REQUIRED PATH}
\NormalTok{tool validation/context completion                        ALLOWED}
\NormalTok{tool semantic approval/equivalence                        REJECTED}
\NormalTok{new BAE family                                             NOT FORCED}
\NormalTok{BA1 / BA2 / BA3 / BA4 reopen                              NOT TRIGGERED}
\NormalTok{BA5                                                        STARTED / NOT CLOSED}
\end{Highlighting}
\end{Shaded}

\subsection{18. Falsification rules}\label{falsification-rules}

Revise this candidate if concrete corpus/integration evidence
demonstrates that:

\begin{enumerate}
\def\labelenumi{\arabic{enumi}.}
\tightlist
\item
  a normative alias must coexist with a canonical token in the same
  semantic domain to preserve necessary shared meaning;
\item
  role semantics cannot be resolved through operator-scoped canonical
  keys;
\item
  a flat registry is actually necessary for interoperability and domain
  separation loses required semantics;
\item
  a semantic-kind/controlled-value extension cannot be governed without
  turning the tool into authority;
\item
  immutable registry revisions are insufficient to replay the meaning
  used by an accepted BA baseline;
\item
  method projection semantics require method-specific taxonomy to enter
  the shared BA registry; or
\item
  the extension decision tree routes a genuine shared semantic need to
  the wrong layer.
\end{enumerate}

Until such evidence exists, synonym/alias machinery remains unnecessary
complexity.

\subsection{19. Smallest unresolved set after
BA5-T2}\label{smallest-unresolved-set-after-ba5-t2}

T1 and T2 now cover named referents, BA2 operators, operator-scoped
roles, semantic kinds, controlled values, exact-token enforcement,
extension governance and the tool authority boundary.

The smallest remaining BA5 task is an integrated closure review that
tries to break the combined contract across:

\begin{itemize}
\tightlist
\item
  referent rename plus registry revision change;
\item
  human and incompatible method projections;
\item
  current facial/order examples;
\item
  extensible versus closed registry domains;
\item
  alias pressure; and
\item
  tool completion/extension workflow.
\end{itemize}

If no material counterexample survives, BA5 may close for current thesis
scope. Do not start BA6 before that closure is official.

\end{document}
